\scanpage{9}
The previous constructions rely heavily on the existence and choice of a field of representatives. Even this being granted, the question of caracterising $Q$ as an actual \uline{scheme} \inserted{over $k$}, i.e. as a functor on the category of $k$-algebras, has not been touched. In fact, the description I can give of this functor is \inserted{(tautological, and\deleted{)} is )} in terms of a particular choice of a resolution $X$, and it really gives \deleted{rather} a description pf $P$ rather than $Q$. \deleted{To} One then still has to define $Q$ as the cokernel of the silly morphism $\underline{D}\handwritten{\longrightarrow}P$. As a matter of fact, even if we restrict to argument\unsure{s} $k'$ which are, say, finite extensions of $k$ (but not necessarily \unsure{separable} ones), I cannot describe $Q(k')$ in terms of local Picard groups of local rings as $A'=A\handwritten{\otimes}_k k'$. The trouble, of course, is that $X\handwritten{\otimes}_k k'$ will in general no longer be a resolution of $S'$, i.e. it will no longer be regular. However, such geometric descriptions will be possible provided we take as arguments algebras which are \deleted{smooth of formally smooth}\inserted{\handwritten{geometrically regular}} over $k$, so \deleted{\unsure{taht}} that tensoring with them will not destroy the regularity of $X$. (Thus, the case when $k'$ is a power series ring, for instance the completion of the local ring of $Q$ at the origin, $Q$ being smooth, is of interest ...) Let's now try to be specific. By \unsure{definition} \inserted{\handwritten{more or less,}} for any \inserted{\handwritten{affine}} scheme $T$ over $k$, we have
\margin{\figure{\handwritten{\begin{tikzpicture}[x=.56cm,y=.56cm]
\draw[line width=.4pt] (-1.7,1.55)--(-.1,2.3) .. controls (.7,2.65) and (1.55,1.95) .. (1.55,.7) .. controls (1.65,-.6) and (.85,-1.15) .. (-.15,-1.55)--(-1.7,-2.2);
\node[inner sep=0pt,align=left,rotate=14] at (-.05,.15) {Affineness of\\$T$ used to\\have $\begin{aligned}H^1(T,\underline{O}_T)&={}\\H^2(T,\underline{O}_T)&=0\end{aligned}$};
\end{tikzpicture}}}}
\[
P(T)=\underset{\handwritten{\longleftarrow}}{\lim}P_n(T)\handwritten{\hookleftarrow}\underset{\handwritten{n}}{\underset{\handwritten{\longleftarrow}}{\lim}}\operatorname{Pic}(X_{n_T}\deleted{\unsure{,\underline{O}_{X_{nT}}}})\handwritten{/\operatorname{Pic}(T)}\qquad (X_{n\,T}=X_n\mathrm{x}_{\unsure{\mathbb{k}|k}} T)\ .
\tag*{\handwritten{(18)}}
\]
In fact, it is not hard to check that the left hand side functor is the étale sheaf (on $(\mathrm{Sch})_{/k}$): associated to the \deleted{\unsure{left hand}} right-hand side functor. This comes from the more precise statement that the obstruction, for an element of $P(T)$, to come from an element of $\deleted{\operatorname{Pic}(X_{\unsure{nT|2T}})}\underset{\handwritten{n}}{\underset{\handwritten{\longleftarrow}}{\lim}}\operatorname{Pic}(X_{nT})$, lies in $H^2(T_{\text{ét}},\underline{G}_m)$ \handwritten{(} at least for $T$ affine, \handwritten{\unsure{(}}which is used to insure that it's $H^{\handwritten{2}}$ and $H^{\handwritten{3}}$ with values in $\underline{O}_T$ vanish\handwritten{\unsure{)}}, For simplicity, I have assumed here that $A$ is normal\handwritten{)}. As we have seen
% archive sheet number: 8
